\documentclass[11pt,twocolumn]{article}
\usepackage{amsmath}
\title{On the Fidelity of Document Conversion}\author{A. Researcher \and B. Scientist}\date{October 2026}
\begin{document}\maketitle
\begin{abstract}We study how documents survive conversion between page-description and word-processing formats, and report measurements on a varied corpus.\end{abstract}
\section{Introduction}
Converting a page into an editable document is a classic problem\footnote{See the extensive literature on layout analysis.}. The difficulty lies in recovering paragraphs, columns and tables from positioned glyphs. Hyphenation, ligatures such as ``fi'' and ``fl'', and kerning further complicate matters. \textbf{Bold text}, \emph{emphasis} and \texttt{monospace} must be preserved.
\section{Method}
We render each page and compare baselines. The error is
\begin{equation} E = \frac{1}{N}\sum_{i=1}^{N} |y_i - \hat{y}_i| \end{equation}
where $y_i$ is the original baseline. Table~\ref{t} lists the results.
\begin{table}[h]\centering\begin{tabular}{|l|r|r|}\hline Set & Pages & Error \\\hline A & 12 & 1.2 \\ B & 40 & 0.8 \\\hline\end{tabular}\caption{Results.}\label{t}\end{table}
\section{Discussion}
Glyph column ligature table corpus measurement corpus editable font table corpus layout editable fidelity layout measurement ligature kerning table document layout layout layout layout editable font fidelity layout renderer kerning measurement corpus kerning conversion kerning kerning measurement hyphen layout fidelity table baseline hyphen table document renderer fidelity renderer font hyphen hyphen corpus renderer editable paragraph corpus kerning editable fidelity baseline conversion conversion column measurement renderer table baseline renderer editable conversion corpus layout corpus paragraph hyphen editable baseline baseline renderer kerning layout font kerning editable renderer conversion conversion measurement ligature layout editable renderer glyph renderer font fidelity paragraph corpus conversion font renderer fidelity corpus conversion fidelity conversion layout document measurement layout kerning baseline baseline column ligature paragraph column column layout measurement.

Layout ligature kerning ligature table baseline conversion hyphen column baseline baseline ligature renderer baseline ligature hyphen measurement document corpus corpus table layout hyphen editable document fidelity font ligature table ligature renderer font fidelity layout kerning layout editable glyph paragraph baseline measurement renderer fidelity kerning renderer measurement kerning renderer layout editable document fidelity paragraph hyphen glyph font paragraph hyphen column column hyphen hyphen baseline fidelity ligature glyph layout paragraph font measurement baseline renderer paragraph editable font conversion table font fidelity font corpus table editable hyphen renderer corpus layout document editable hyphen layout baseline font document glyph document fidelity font ligature table editable conversion corpus kerning column paragraph column glyph baseline baseline font ligature document renderer ligature conversion document document table hyphen.

Kerning corpus glyph table document paragraph fidelity column editable glyph glyph document table editable column kerning column ligature conversion hyphen table measurement ligature table paragraph hyphen layout layout column fidelity table paragraph font kerning fidelity baseline table measurement baseline kerning baseline table fidelity editable hyphen ligature corpus document table font document paragraph layout layout hyphen document measurement editable document editable column column document measurement table ligature font corpus conversion ligature baseline font hyphen font kerning conversion column ligature column measurement column document kerning editable hyphen paragraph document baseline document hyphen kerning document table column kerning kerning layout kerning editable column ligature column column layout layout hyphen conversion corpus corpus glyph table renderer document column renderer baseline baseline glyph glyph document.

Hyphen table renderer hyphen glyph font glyph paragraph document font baseline hyphen fidelity baseline paragraph kerning ligature column measurement fidelity ligature measurement measurement layout editable document baseline ligature corpus layout fidelity layout paragraph conversion glyph glyph glyph ligature ligature editable editable baseline column kerning corpus layout baseline renderer document renderer measurement kerning kerning document corpus corpus kerning fidelity document ligature kerning paragraph column renderer conversion baseline renderer font hyphen hyphen hyphen conversion baseline measurement column table renderer editable baseline glyph ligature fidelity font paragraph corpus editable conversion editable renderer baseline paragraph renderer column ligature table ligature column glyph column measurement kerning editable fidelity editable baseline document measurement glyph corpus font table fidelity fidelity table hyphen ligature kerning editable layout font.

Renderer measurement layout layout kerning ligature font baseline hyphen glyph font ligature hyphen ligature measurement baseline conversion corpus fidelity table font editable font hyphen table layout table layout hyphen glyph column renderer conversion hyphen fidelity renderer conversion renderer document layout table measurement measurement conversion hyphen editable document corpus table editable editable font layout ligature renderer font measurement renderer fidelity hyphen baseline measurement renderer font conversion renderer layout editable fidelity editable document column corpus kerning hyphen layout fidelity glyph editable ligature baseline column layout conversion ligature fidelity hyphen glyph measurement ligature corpus baseline measurement renderer paragraph ligature renderer table fidelity column conversion column measurement layout baseline renderer baseline column editable ligature hyphen font renderer font kerning document ligature column column renderer.

Conversion measurement renderer paragraph baseline hyphen ligature conversion kerning editable editable baseline corpus ligature document kerning ligature kerning layout editable document fidelity kerning ligature font column baseline measurement glyph ligature measurement renderer baseline glyph glyph measurement conversion hyphen editable kerning table font hyphen column table kerning editable document corpus table baseline paragraph paragraph layout font paragraph corpus renderer measurement document ligature table baseline table kerning editable kerning corpus measurement editable baseline kerning kerning hyphen measurement editable font measurement ligature document corpus table font column paragraph layout layout corpus document editable hyphen font editable baseline glyph layout layout editable glyph paragraph editable ligature glyph column measurement hyphen layout paragraph paragraph renderer glyph paragraph ligature table fidelity column font layout corpus glyph.

Ligature font measurement editable document ligature ligature kerning kerning paragraph baseline conversion fidelity renderer paragraph conversion fidelity font fidelity column ligature column ligature baseline table glyph paragraph font fidelity paragraph paragraph column renderer corpus renderer conversion table document paragraph glyph paragraph measurement glyph editable measurement layout renderer ligature column ligature document column hyphen paragraph editable paragraph ligature document glyph ligature editable table hyphen table fidelity kerning renderer font document document renderer editable corpus table glyph measurement renderer renderer layout hyphen baseline font conversion editable renderer document table fidelity conversion glyph column paragraph hyphen document fidelity hyphen document conversion ligature document renderer renderer layout renderer table glyph document document document column measurement ligature corpus measurement conversion editable column paragraph glyph paragraph.

Renderer corpus ligature kerning document conversion conversion editable hyphen measurement document renderer baseline layout glyph ligature kerning glyph table baseline fidelity paragraph table ligature table font ligature column renderer column column font baseline renderer fidelity layout conversion corpus hyphen kerning font corpus kerning fidelity measurement conversion font corpus column ligature fidelity font layout editable renderer corpus column editable renderer fidelity paragraph conversion measurement layout font hyphen layout table hyphen renderer document hyphen renderer fidelity renderer fidelity hyphen measurement hyphen glyph renderer measurement glyph baseline ligature layout fidelity paragraph conversion fidelity editable hyphen layout column column layout editable ligature measurement ligature conversion corpus document editable measurement table corpus conversion glyph fidelity glyph layout baseline ligature conversion glyph hyphen fidelity ligature renderer.
\end{document}
